\label{chap:instrumentacion_metodologia}

En este capítulo se describe la instrumentación empleada para adquirir
y procesar las señales acústicas, así como los procedimientos utilizados
para caracterizar experimentalmente los observables desarrollados en
los capítulos anteriores. Se presentan la arquitectura general de
WALO, la configuración de adquisición y los montajes experimentales
para la estimación de distancia, dirección de llegada y movimiento
de una fuente sonora.

La metodología presenta la configuración instrumental y los
procedimientos de los experimentos de distancia y dirección.
También describe el montaje cinemático, la adquisición continua
y la referencia óptica utilizada para contrastar la estimación
Doppler mediante seguimiento en Tracker.

La implementación de hardware y software se encuentra documentada
en repositorios versionados. El diseño del overlay para la plataforma
PYNQ-Z2 corresponde a la versión 1.5.3 del repositorio de hardware
\cite{SiririHW2026}, mientras que las rutinas de adquisición y
procesamiento se encuentran en el repositorio de software
\cite{SiririSW2026}. Estas referencias permiten identificar las
implementaciones utilizadas y consultar su evolución.

\section{Arquitectura e instrumentación de WALO}
\label{sec:arquitectura_walo}

WALO es una plataforma de adquisición y procesamiento de señales
acústicas implementada sobre una tarjeta PYNQ-Z2. El sistema integra
dos canales de recepción analógica, conversión analógico-digital,
procesamiento digital programable en FPGA y una interfaz de control
desarrollada en Python.

La arquitectura se organiza en tres etapas funcionales:

\begin{enumerate}
    \item \textbf{Transducción:} conversión de las variaciones de
    presión acústica en señales eléctricas mediante dos módulos
    de micrófono.

    \item \textbf{Adquisición y procesamiento digital:} conversión
    de las señales analógicas mediante el XADC, acondicionamiento
    digital y transferencia de los datos hacia memoria.

    \item \textbf{Control y análisis:} configuración del hardware,
    captura de registros, extracción de observables y estimación
    de magnitudes físicas mediante rutinas de Python.
\end{enumerate}

\subsection{Unidad de procesamiento y adquisición}

La plataforma utiliza una tarjeta PYNQ-Z2 basada en el dispositivo
Zynq-7000 XC7Z020. Este integra un sistema de procesamiento
ARM y una matriz de lógica programable FPGA. La separación entre
ambos dominios permite ejecutar las tareas de control y análisis
en software, mientras que la adquisición y determinadas operaciones
de procesamiento se realizan mediante bloques de hardware
configurados en la lógica programable.

La conversión analógico-digital se realiza mediante el XADC
integrado en el dispositivo. Los dos canales de recepción se
conectan a las entradas analógicas A0 y A1 de la tarjeta,
correspondientes a los canales auxiliares Vaux1 y Vaux9,
respectivamente.

La arquitectura de adquisición utiliza el secuenciador dual del
XADC para registrar ambos canales de manera simultánea. La
configuración nominal del convertidor corresponde a una tasa
agregada de \SI{1}{MSPS}, equivalente a
\SI{500}{kSPS} por canal antes del diezmado digital. La
sincronización intercanal se verificó experimentalmente mediante
una prueba con una señal común aplicada a ambas entradas. La
descripción del procedimiento y su evidencia técnica se encuentran
en los repositorios de hardware y software anexos
\cite{SiririHW2026,SiririSW2026}.

\subsection{Transductores y excitación acústica}

La recepción acústica se realiza mediante dos módulos de micrófono
electret con preamplificación MAX4466. Cada módulo entrega una
señal analógica cuya amplitud depende de la presión acústica
incidente, de la respuesta del transductor y de la ganancia
configurada en el circuito.

Durante los experimentos de caracterización se utilizó la
configuración de ganancia máxima disponible en ambos módulos.
No se dispone de una calibración absoluta de dicha ganancia ni
de una caracterización independiente de su linealidad. Por esta
razón, las amplitudes registradas se tratan como magnitudes
eléctricas de los canales y no como mediciones directas de
presión acústica en pascales.

La excitación acústica se produjo mediante un buzzer activo
con frecuencia nominal cercana a \SI{2600}{\hertz}. Para los
experimentos pulsados, el buzzer se accionó mediante una salida
digital de la tarjeta, conectada a un transistor 2N2222A que
controla su alimentación. Esta etapa permite generar pulsos
acústicos bajo el control temporal del sistema digital.

\subsection{Arquitectura de procesamiento digital}

La lógica programable ejecuta la adquisición de los canales,
la reducción programable de la tasa de muestreo y la transferencia
de los registros mediante interfaces DMA. La arquitectura de
hardware utilizada en los experimentos corresponde a la versión
1.5.3 del overlay.

En los experimentos de distancia y dirección se utilizó un
factor de diezmado $M=10$, correspondiente a una tasa efectiva
de \SI{50}{kSPS} por canal en los registros procesados. En
consecuencia, debe distinguirse esta tasa de salida de la tasa
nominal de conversión del XADC, de \SI{500}{kSPS} por canal
antes del diezmado.

\begin{figure}[!htbp]
    \centering
    
\begin{tikzpicture}[
    >=Latex,
    font=\small,
    block/.style={
        draw,
        rounded corners=2pt,
        align=center,
        minimum width=2.5cm,
        minimum height=0.9cm,
        inner sep=5pt
    },
    smallblock/.style={
        draw,
        rounded corners=2pt,
        align=center,
        minimum width=2.3cm,
        minimum height=0.8cm,
        inner sep=5pt
    },
    arrow/.style={->, thick}
]

% Recepción analógica
\node[block] (mic1) at (0,2.0) {Micrófono 1};
\node[block] (mic2) at (0,0.5) {Micrófono 2};

\node[block] (xadc) at (3.4,1.25)
    {XADC\\Conversión A/D};

\node[block] (zynq) at (6.9,1.25)
    {Zynq-7000\\PYNQ-Z2};

\node[smallblock] (dma) at (10.2,2.15)
    {DMA\\Transferencia};

\node[smallblock] (processing) at (10.2,0.35)
    {Procesamiento\\digital};

\node[block] (software) at (13.6,2.15)
    {Software\\de control};

\node[smallblock] (analysis) at (13.6,0.35)
    {Análisis\\de señales};

% Conexiones de recepción
\draw[arrow] (mic1.east) -- (xadc.west);
\draw[arrow] (mic2.east) -- (xadc.west);

\draw[arrow] (xadc.east) -- (zynq.west);
\draw[arrow] (zynq.east) -- (dma.west);
\draw[arrow] (zynq.east) -- (processing.west);

\draw[arrow] (dma.east) -- (software.west);
\draw[arrow] (processing.east) -- (software.south west);
\draw[arrow] (software.south) -- (analysis.north);

% Ruta de excitación desde la tarjeta
\node[smallblock] (pulse) at (3.4,-1.25)
    {Pulso digital\\de salida AR2};

\node[smallblock] (driver) at (7.1,-1.25)
    {Etapa de excitación\\Transistor 2N2222A};

\node[block] (buzzer) at (10.6,-1.25)
    {Buzzer activo};

\draw[arrow] (pulse.east) -- (driver.west);
\draw[arrow] (driver.east) -- (buzzer.west);

% Relación de la salida digital con la tarjeta
\draw[arrow, dashed] (zynq.south) -- ++(0,-0.55)
    -| (pulse.north);

\end{tikzpicture}

    \caption{Arquitectura funcional de WALO. Se distinguen la
    recepción analógica, la adquisición mediante XADC, el
    procesamiento en FPGA y la transferencia de datos hacia
    el sistema de procesamiento.}
    \label{fig:arquitectura_walo}
\end{figure}

\section{Configuración de adquisición y procesamiento}
\label{sec:configuracion_adquisicion}

La configuración de adquisición depende del observable que se
desea estimar. Los experimentos de distancia y dirección emplean
excitación pulsada, mientras que el experimento cinemático requiere
una adquisición continua para seguir la variación temporal de
la frecuencia recibida.

\subsection{Adquisición pulsada}

En los experimentos de distancia y dirección, el sistema genera
un pulso de excitación mediante la salida digital AR2. El instante
de activación proporciona una referencia temporal para la
captura de la respuesta acústica.

La adquisición se realiza simultáneamente en ambos canales.
El registro temporal permite identificar la llegada de la señal
directa y extraer las magnitudes asociadas a la ventana de
observación. La duración de la ventana se define en relación
con la frecuencia nominal del buzzer y con el número de ciclos
analizados.

Para la caracterización de distancia se utilizó una ventana
correspondiente a tres ciclos de la señal nominal. Esta
elección se relacionó con la geometría del montaje: la mesa,
situada aproximadamente \SI{50}{\centi\metre} por debajo
de la fuente y del micrófono, constituía la superficie
reflectante más cercana. La duración de la ventana es

\begin{equation}
    T_{\mathrm{gate}} =
    \frac{K}{f_0},
    \label{eq:ventana_gate_exp}
\end{equation}

donde $K=3$ y $f_0$ es la frecuencia nominal de la fuente.
Para $f_0=\SI{2600}{\hertz}$, la duración es aproximadamente
\SI{1.15}{\milli\second}.

La ventana se seleccionó para restringir el análisis a la
respuesta inicial y reducir la contribución de reflexiones
posteriores. La mesa es la superficie reflectante más cercana;
sin embargo, no se dispone aquí de una determinación completa
de las longitudes de las trayectorias directa y reflejada que
permita demostrar que la primera reflexión llega después del
intervalo de \SI{1.15}{\milli\second}. Por ello, la ventana
privilegia la respuesta inicial, pero no se considera una
garantía de exclusión de todas las reflexiones. El recinto no
se trata como anecoico.

\subsection{Adquisición continua}

Para el estudio del movimiento de la fuente se registraron
continuamente ambos canales durante diez ensayos de
\SI{10}{\second} cada uno. La tasa de adquisición almacenada
fue de \SI{50}{kSPS} por canal, después de aplicar el factor
de diezmado $M=10$. Cada ensayo contiene, por tanto,
500\,000 muestras por canal. La excitación del buzzer se
mantuvo activa durante todo el registro y la adquisición
comenzó antes de iniciar el desplazamiento, de modo que
incluyera un intervalo inicial de reposo.

A partir de las señales se estimó la frecuencia recibida en
ventanas sucesivas. Para la caracterización cinemática, la
secuencia de estimaciones se procesó a una tasa temporal de
\SI{100}{Hz}. Esta tasa corresponde a la salida del
procesamiento cinemático y no a la tasa de muestreo del
convertidor. Las dos frecuencias se utilizan para construir
estimaciones de velocidad relativa conforme al modelo Doppler
desarrollado en el capítulo anterior.

\subsection{Calibración temporal de los canales}

La estimación absoluta de distancia mediante ToA requiere
caracterizar el retardo entre la orden de emisión y la
detección del arribo acústico. Este retardo incluye
contribuciones de la etapa de excitación, del transductor
emisor, de los micrófonos, de la electrónica de recepción
y del procesamiento digital.

Para cada canal se utilizó una medición de referencia a
\SI{10}{\centi\metre}. Si el tiempo de llegada medido a
esa distancia es $t_{\mathrm{arr,ref}}$, el retardo efectivo
se determina mediante

\begin{equation}
    t_{\mathrm{off}} =
    t_{\mathrm{arr,ref}}
    -t_{\mathrm{em}}
    -\frac{r_{\mathrm{ref}}}{c(T)},
    \label{eq:offset_calibracion_exp}
\end{equation}

donde $t_{\mathrm{em}}$ representa el instante de referencia
asociado a la orden digital de excitación, y $c(T)$ la
velocidad del sonido utilizada para la temperatura del ensayo.
Esta referencia corresponde al evento eléctrico de activación
del buzzer y no debe confundirse con el instante exacto en
que comienza la radiación acústica. El retardo entre ambos
eventos queda incluido en el retardo efectivo determinado
mediante la calibración.

La calibración se realizó de manera independiente para
cada micrófono. Los valores obtenidos se almacenaron en el
perfil de calibración utilizado por las rutinas de
estimación. El retardo efectivo reúne las contribuciones
de la cadena de excitación, recepción y procesamiento bajo
la configuración empleada.

\section{Metodología de caracterización experimental}
\label{sec:metodologia_experimental}

\subsection{Caracterización experimental de distancia}
\label{sec:exp01_distancia}

El primer experimento caracteriza la respuesta de cada
canal de recepción en función de la distancia entre el
buzzer y el micrófono. Se estudian las variaciones de
amplitud, energía digital y tiempo de llegada, con el
propósito de establecer modelos empíricos para la
estimación radial.

\subsubsection{Montaje experimental}

El montaje utiliza un buzzer activo como fuente acústica
y un micrófono como receptor. Los dos canales se
caracterizan de manera independiente, manteniendo
la misma configuración geométrica y electrónica.

La fuente y el micrófono se ubicaron a una misma altura
aproximada de \SI{50}{\centi\metre} sobre una mesa de
madera. La orientación fue axial: la salida del buzzer
se dirigió hacia la cara del micrófono. La distancia
se midió entre el centro de la salida acústica del
buzzer y el centro de la membrana del micrófono.

Las distancias se determinaron mediante una cinta métrica 
con resolución indicada de \SI{1}{\milli\metre}.
La referencia geométrica también depende de la alineación
y de la identificación de los centros de los transductores.

\begin{figure}[!htbp]
    \centering
    
\begin{tikzpicture}[
    >=Latex,
    font=\small,
    source/.style={
        circle,
        draw=orange!75!black,
        fill=orange!25,
        thick,
        minimum size=6mm,
        inner sep=0pt
    },
    receiver/.style={
        circle,
        draw=gray!65,
        fill=gray!20,
        thick,
        minimum size=6mm,
        inner sep=0pt
    },
    direct/.style={->, thick, green!45!black},
    reflected/.style={->, thick, dashed, red!75!black},
    dimension/.style={<->, thick}
]

% Vista superior
\node[font=\bfseries] at (-3.0,3.25) {Vista superior};

\node[source] (sourceTop) at (-3.0,1.85) {};
\node[above=5pt] at (sourceTop.north) {Buzzer};

\node[receiver] (receiverTop) at (-3.0,-0.75) {};
\node[below=7pt] at (receiverTop.south) {Micrófono};

\draw[direct] (sourceTop.south) -- (receiverTop.north)
    node[midway,right=7pt,align=left] {Eje de\\propagación};

\draw[dimension] (-4.0,1.85) -- (-4.0,-0.75)
    node[midway,left] {$r$};

% Vista lateral
\node[font=\bfseries] at (3.7,3.25) {Vista lateral};

% Mesa con espesor
\filldraw[
    fill=gray!12,
    draw=black,
    thick
] (0,-0.45) rectangle (7.5,-0.18);

\node[font=\scriptsize] at (3.75,-0.315) {Mesa};

% Soportes
\draw[thick] (1.0,-0.18) -- (1.0,1.05);
\draw[thick] (6.3,-0.18) -- (6.3,1.05);

% Micrófono a la izquierda
\node[receiver] (receiverSide) at (1.0,1.45) {};
\node[above=6pt] at (receiverSide.north) {Micrófono};

% Buzzer a la derecha
\node[source] (sourceSide) at (6.3,1.45) {};
\node[above=6pt] at (sourceSide.north) {Buzzer};

% Trayectoria directa
\draw[direct]
    (sourceSide.west) -- (receiverSide.east)
    node[midway,above=7pt] {$r$};

% Trayectoria reflejada
\coordinate (bounce) at (3.65,-0.12);

\draw[reflected]
    (sourceSide.south west) -- (bounce)
    -- (receiverSide.south east);

\node[
    font=\scriptsize,
    text=red!75!black,
    fill=white,
    inner sep=2pt
] at (3.65,0.22) {Reflexión en mesa};

% Altura
\draw[dimension] (7.0,-0.18) -- (7.0,1.05)
    node[midway,right=5pt] {\SI{50}{cm}};

% Recuadro inferior, separado de la mesa
\node[
    draw,
    rounded corners=2pt,
    fill=gray!5,
    text width=6.5cm,
    align=center,
    inner sep=7pt,
    font=\scriptsize
] at (3.75,-1.35) {
    Fuente y receptor a altura aproximadamente constante,\\
    enfrentados sobre el eje de propagación
};

\end{tikzpicture}

    \caption{Montaje experimental para la caracterización
    radial. La fuente y el micrófono se mantienen a una
    altura aproximadamente constante y enfrentados sobre
    el eje de propagación.}
    \label{fig:montaje_exp01}
\end{figure}

\subsubsection{Condiciones de medición}

La fuente se operó a una frecuencia nominal cercana a
\SI{2600}{\hertz}. La temperatura de referencia utilizada
en el procesamiento fue de \SI{20}{\degreeCelsius},
con una velocidad del sonido aproximada de
\SI{343.2}{\metre\per\second}.

Se realizaron mediciones en nueve posiciones radiales:

\begin{equation}
    r \in
    \{15,20,25,30,40,50,60,80,100\}
    \ \si{\centi\metre}.
\end{equation}

En cada posición se registraron veinte pulsos por canal.
Por tanto, cada micrófono cuenta con ciento ochenta
registros distribuidos en las nueve distancias.

La configuración de ganancia de ambos módulos se mantuvo
en el máximo disponible. La caracterización se realizó
primero para un canal y luego para el otro, conservando
la geometría del montaje.

\subsubsection{Magnitudes registradas}

Para cada pulso se conservaron las señales temporales
y las magnitudes calculadas por la rutina de adquisición.
Entre los observables físicos de interés se encuentran:

\begin{itemize}
    \item Amplitud RMS de la señal recibida;
    \item Magnitud digital asociada a la energía directa;
    \item Tiempo de llegada detectado;
\end{itemize}

La amplitud RMS y la energía directa son observables
instrumentales obtenidos a partir de las señales eléctricas
digitalizadas. La amplitud RMS caracteriza el nivel eléctrico
de cada canal, mientras que $E_{\mathrm{direct}}$ es una
magnitud digital dependiente de la implementación, de la
ventana temporal, de la tasa de muestreo y de la representación
numérica. Ninguna de las dos constituye, con la configuración
utilizada, una medición absoluta de presión acústica en
pascales o de energía acústica en joules.

\subsubsection{Calibración y caracterización radial}

Los parámetros de los modelos radiales se determinan
a partir de las mediciones de amplitud y energía digital.
Para la amplitud se considera el modelo

\begin{equation}
    A(r)=\frac{k_A}{r},
\end{equation}

mientras que para la magnitud digital asociada a la
energía se considera

\begin{equation}
    \mathcal{E}_{\mathrm{dig}}(r)=\frac{k_E}{r^2}.
\end{equation}

Los parámetros $k_A$ y $k_E$ se estiman de manera
independiente para cada canal. Estos coeficientes
representan calibraciones empíricas del sistema
completo bajo las condiciones del montaje.

La calibración temporal por canal se efectuó a una
distancia de referencia de \SI{10}{\centi\metre}.
El retardo efectivo obtenido se incorporó en la
estimación de distancia por ToA.

La evaluación de los modelos radiales debe distinguir
entre los datos utilizados para determinar los
coeficientes y aquellos empleados para evaluar su
capacidad predictiva. Cuando se utilizan las mismas
mediciones para ambas operaciones, el error obtenido
describe el ajuste dentro de la muestra y no constituye
una validación independiente.

\subsubsection{Condiciones de propagación}

El montaje se realizó en un recinto interior con
superficies reflectantes. La elevación de los
transductores sobre la mesa y el uso de excitación
pulsada buscan reducir la contribución de reflexiones
tardías dentro de la ventana temporal analizada.

La ventana de observación se selecciona para analizar
la respuesta inicial y privilegiar la contribución de
la trayectoria directa. Esta elección puede reducir
la influencia de reflexiones tardías, pero no garantiza
que todas las contribuciones reflejadas queden fuera
de la ventana analizada.

La presencia de superficies reflectantes, incluida la
mesa de madera, impide asumir condiciones anecoicas.
Por tanto, la compatibilidad de los datos con las
dependencias radiales esperadas se evalúa
experimentalmente y se interpreta dentro de las
condiciones concretas del montaje.

\subsection{Caracterización experimental de dirección de llegada}
\label{sec:exp02_angular}

El segundo experimento evalúa la capacidad del arreglo
de dos micrófonos para estimar la dirección de llegada
de una señal acústica. Se comparan dos procedimientos:
la estimación basada en diferencias de tiempos de
llegada (TDOA) y la estimación basada en diferencias
de fase intercanal (PDoA).

\subsubsection{Geometría del arreglo}

Los dos micrófonos se dispusieron sobre una línea
recta, con una separación entre centros de
\SI{5}{\centi\metre}. El centro geométrico del arreglo
se utilizó como referencia para definir la posición
angular de la fuente.

El buzzer se ubicó a una distancia radial aproximada
de \SI{30}{\centi\metre} del centro del arreglo y se
orientó hacia ambos micrófonos. La altura de los
transductores se mantuvo aproximadamente en
\SI{50}{\centi\metre} sobre la mesa.

\begin{figure}[!htbp]
    \centering
    
% Montaje experimental para la caracterización angular.
% El entorno figure, caption y label se definen en el capítulo.

\begin{tikzpicture}[
    >=Stealth,
    line cap=round,
    line join=round,
    font=\small,
    x=1cm,
    y=1cm
]

    % =================================================
    % 1. GEOMETRIA DEL ARREGLO
    % =================================================

    % Centro del arreglo desplazado hacia la izquierda
    \coordinate (O) at (-3,0);

    % Micrófonos: mayor separación horizontal
    \coordinate (Mone) at (-3.425,0);
    \coordinate (Mtwo) at (-2.575,0);

    % Eje horizontal principal
    \draw[line width=0.95pt]
        (-5.35,0) -- (-0.65,0);

    % Eje vertical principal
    \draw[line width=0.95pt]
        (-3,-0.05) -- (-3,4.35);

    % =================================================
    % 2. POSICIONES ANGULARES
    % =================================================

    \foreach \ang in {-60,-45,-30,-15,15,30,45,60}{
        \draw[
            densely dashed,
            line width=0.8pt
        ]
        ({-3+4.0*sin(\ang)},{4.0*cos(\ang)})
        --
        ({-3+0.38*sin(\ang)},{0.38*cos(\ang)});
    }

    % =================================================
    % 3. FUENTE ACUSTICA
    % =================================================

    % Posición ilustrativa de la fuente a +30 grados
    \coordinate (S) at (
        {-3+3.55*sin(30)},
        {3.55*cos(30)}
    );

    \node[
        circle,
        draw=orange!85!black,
        fill=orange!18,
        line width=1pt,
        minimum size=6mm,
        inner sep=0pt
    ] at (S) {};

    \node[
        anchor=south,
        inner sep=2pt
    ] at (-1.22,3.48) {Fuente};

    % =================================================
    % 4. REFERENCIA ANGULAR
    % =================================================

    \draw[
        line width=0.85pt
    ]
    (-3,0.78)
    arc[
        start angle=90,
        end angle=60,
        radius=0.78
    ];

    \node[
        fill=white,
        inner sep=2pt
    ] at (-2.77,0.83) {$\theta$};

    % =================================================
    % 5. DISTANCIA RADIAL
    % =================================================

    \node[
        fill=white,
        inner sep=2pt
    ] at (-4.05,2.35)
    {$r=30\,\mathrm{cm}$};

    % =================================================
    % 6. MICROFONOS
    % =================================================

    % Micrófono 1
    \node[
        circle,
        draw=black!45,
        fill=black!12,
        line width=0.9pt,
        minimum size=5.8mm,
        inner sep=0pt
    ] at (Mone) {};

    % Micrófono 2
    \node[
        circle,
        draw=black!45,
        fill=black!12,
        line width=0.9pt,
        minimum size=5.8mm,
        inner sep=0pt
    ] at (Mtwo) {};

    % Identificación individual con líneas guía
    \draw[black!65,line width=0.55pt]
        (-3.425,-0.15) -- (-3.72,-0.48);

    \draw[black!65,line width=0.55pt]
        (-2.575,-0.15) -- (-2.28,-0.48);

    \node[
        anchor=north
    ] at (-3.76,-0.48) {Mic. 1};

    \node[
        anchor=north
    ] at (-2.24,-0.48) {Mic. 2};

    % =================================================
    % 7. DISTANCIA ENTRE MICROFONOS
    % =================================================

    \draw[densely dashed,black!45]
        (-3.425,-0.22) -- (-3.425,-1.12);

    \draw[densely dashed,black!45]
        (-2.575,-0.22) -- (-2.575,-1.12);

    % Cota de separación
    \draw[
        <->,
        line width=0.75pt
    ]
        (-3.425,-1.02) -- (-2.575,-1.02);

    \node[
        fill=white,
        inner sep=2pt
    ] at (-3,-1.30)
    {$d=5\,\mathrm{cm}$};

    % Centro geométrico
    \fill (O) circle (1.4pt);

    % =================================================
    % 8. CONFIGURACION EXPERIMENTAL
    % =================================================

    \node[
        draw=black!65,
        rounded corners=2pt,
        fill=black!2,
        line width=0.65pt,
        text width=4.15cm,
        align=left,
        inner sep=9pt,
        anchor=west
    ] at (1.05,1.35) {%
        \textbf{Configuración experimental}\\[6pt]
        Línea base: $5\,\mathrm{cm}$\\[3pt]
        Distancia nominal: $30\,\mathrm{cm}$\\[3pt]
        Ángulos: $-60^\circ$ a $+60^\circ$\\[3pt]
        Paso angular: $15^\circ$
    };

\end{tikzpicture}

    \caption{Montaje del experimento angular. Los micrófonos
    forman una línea base de \SI{5}{\centi\metre} y la
    fuente se desplaza sobre posiciones angulares
    definidas respecto al centro del arreglo.}
    \label{fig:montaje_exp02}
\end{figure}

\subsubsection{Posiciones angulares y referencia}

La posición de la fuente se estableció mediante
marcas angulares respecto al eje normal al arreglo.
Se consideraron nueve posiciones:

\begin{equation}
    \theta_{\mathrm{ref}}\in
    \{-60^\circ,-45^\circ,-30^\circ,-15^\circ,
    0^\circ,15^\circ,30^\circ,45^\circ,60^\circ\}.
\end{equation}

La distancia radial se mantuvo aproximadamente
constante durante el barrido. La referencia angular
se determinó mediante el montaje geométrico y el
instrumento de medición angular disponible. La
incertidumbre de esta referencia se estima en
aproximadamente \SI{1}{\degree}; este valor debe
entenderse como una estimación del procedimiento
de posicionamiento y no como una resolución angular
intrínseca del sistema WALO.

Para cada posición se registraron diez pulsos, para
un total de noventa registros en el conjunto de datos.

\subsubsection{Adquisición y estimadores evaluados}

La adquisición se realizó con factor de diezmado
$M=10$ y una tasa de \SI{50}{kSPS} por canal. La excitación pulsada
proporciona una referencia temporal común y permite
analizar la respuesta inicial de ambos micrófonos.

Para cada registro se obtuvieron estimaciones
angulares mediante dos métodos:

\begin{enumerate}
    \item \textbf{TDOA:} se identifican los tiempos
    de llegada en ambos canales y se convierte su
    diferencia en una estimación angular mediante
    la geometría de la línea base.

    \item \textbf{PDoA:} se estima la diferencia de
    fase intercanal de la componente acústica de
    interés y se obtiene el ángulo mediante el
    modelo interferométrico.
\end{enumerate}

La evaluación de ambos métodos utiliza el mismo
conjunto de posiciones de referencia. Para cada
posición se calculan primero las estadísticas de
las diez adquisiciones correspondientes. Las métricas
globales se obtienen posteriormente a partir de los
resultados por posición, evitando tratar las noventa
capturas como si correspondieran a una única
condición angular.


\subsubsection{Control de sincronización y coherencia intercanal}

La estimación angular depende de que las diferencias
temporales y de fase entre los canales correspondan
a la propagación acústica y no a retardos introducidos
por la adquisición.

La arquitectura utiliza adquisición simultánea de ambos
canales y la sincronización intercanal fue verificada
experimentalmente mediante una señal común aplicada a
las dos entradas. La descripción de la prueba y su
evidencia se encuentran disponibles en los repositorios
de hardware y software anexados
\cite{SiririHW2026,SiririSW2026}. La caracterización
de la respuesta analógica de cada módulo se considera
por separado al interpretar las diferencias de fase y
de tiempo.

\subsection{Protocolo de caracterización cinemática}
\label{sec:exp03_cinematica}

El tercer experimento está diseñado para estudiar
la estimación de velocidad de una fuente acústica
en movimiento mediante el desplazamiento Doppler
diferencial registrado por dos micrófonos fijos.
A diferencia de los experimentos anteriores, que
utilizan excitación pulsada, esta medición requiere
una adquisición continua durante el movimiento.

El montaje permite estudiar la estimación de velocidad
mediante desplazamiento Doppler y contrastarla con una
referencia óptica independiente obtenida mediante Tracker.
La campaña experimental comprende diez ensayos de
\SI{10}{\second}, con registros acústicos de ambos canales
y seguimiento óptico de la fuente móvil. La comparación
entre ambas mediciones se realiza después de alinear sus
bases temporales; esta alineación se trata como parte del
procesamiento y no como una calibración independiente de
latencia.

\subsubsection{Montaje mecánico}

El montaje utiliza un riel rectilíneo de longitud nominal
aproximada de \SI{2}{\metre}. Un carro se desplaza sobre
el riel y lleva instalado el buzzer, que constituye la
fuente acústica móvil. El movimiento se inicia mediante
un empuje manual y queda limitado por topes mecánicos
rígidos en los extremos; no se utiliza un elemento
elástico para impulsar ni para restaurar el movimiento.

Se disponen dos micrófonos fijos en extremos opuestos
del recorrido. El micrófono 1 se ubica en el extremo
inicial y el micrófono 2 en el extremo final. El buzzer
se orienta hacia el micrófono 1 y en sentido opuesto
al micrófono 2. Aunque el riel tiene una longitud
nominal de \SI{2}{\metre}, las posiciones efectivas
de los extremos del recorrido, referidas al origen
geométrico empleado en la caracterización, son
\SI{9.4(1)}{\centi\metre} y
\SI{175.6(1)}{\centi\metre}. El recorrido útil
resultante es \SI{166.2(14)}{\centi\metre}.
Estas magnitudes describen el intervalo mecánico
utilizable y no deben confundirse con la posición
óptica instantánea del buzzer.

La masa del conjunto móvil no se utiliza en la
estimación Doppler de velocidad y no se asigna un
valor en este capítulo, pues no se dispone de una
medición documentada.

\begin{figure}[!htbp]
    \centering
    
% Montaje experimental cinemático.
% Sin entorno figure, caption ni label.

\begin{tikzpicture}[
    >=Stealth,
    line cap=round,
    line join=round,
    font=\small,
    x=1cm,
    y=1cm
]

% =================================================
% 1. RIEL
% =================================================

\draw[line width=1.2pt] (-0.8,0) -- (10.8,0);

\node[anchor=north] at (5,-0.17)
    {Riel de $2\,\mathrm{m}$};

% =================================================
% 2. BLOQUES DE LOS EXTREMOS
% =================================================

\draw[
    fill=black!16,
    draw=black,
    line width=0.8pt
] (-0.15,0) rectangle (0.15,0.72);

\draw[
    fill=black!16,
    draw=black,
    line width=0.8pt
] (9.85,0) rectangle (10.15,0.72);

% =================================================
% 3. MICROFONOS SOBRE LOS BLOQUES
% =================================================

% Micrófono 1
\draw[line width=1pt,black!65]
    (0,0.72) -- (0,1.62);

\node[
    circle,
    draw=black!45,
    fill=black!12,
    line width=0.9pt,
    minimum size=6mm,
    inner sep=0pt
] at (0,1.62) {};

\node[anchor=south] at (0,1.94)
    {Mic. 1};

% Micrófono 2
\draw[line width=1pt,black!65]
    (10,0.72) -- (10,1.62);

\node[
    circle,
    draw=black!45,
    fill=black!12,
    line width=0.9pt,
    minimum size=6mm,
    inner sep=0pt
] at (10,1.62) {};

\node[anchor=south] at (10,1.94)
    {Mic. 2};

% =================================================
% 4. CAUCHOS / RESORTES DE RETORNO
% =================================================

% Resorte izquierdo
\draw[
    red!75!black,
    line width=1.5pt
]
    (0.15,0.48)
    -- (0.22,0.48)
    -- (0.27,0.57)
    -- (0.33,0.39)
    -- (0.39,0.57)
    -- (0.45,0.39)
    -- (0.51,0.57)
    -- (0.57,0.39)
    -- (0.63,0.48)
    -- (0.70,0.48);

% Resorte derecho
\draw[
    red!75!black,
    line width=1.5pt
]
    (9.30,0.48)
    -- (9.37,0.48)
    -- (9.43,0.57)
    -- (9.49,0.39)
    -- (9.55,0.57)
    -- (9.61,0.39)
    -- (9.67,0.57)
    -- (9.73,0.39)
    -- (9.79,0.48)
    -- (9.85,0.48);

% =================================================
% 5. CARRITO MOVIL
% =================================================

% Cuerpo bajado 0.11 cm
\draw[
    rounded corners=2pt,
    fill=black!4,
    draw=black,
    line width=0.85pt
] (1.35,0.27) rectangle (2.65,0.87);

% Ruedas: ahora su punto inferior coincide con y=0
\draw[
    fill=white,
    draw=black,
    line width=0.7pt
] (1.65,0.105) circle (0.105);

\draw[
    fill=white,
    draw=black,
    line width=0.7pt
] (2.35,0.105) circle (0.105);

\node at (2.00,0.57)
    {Carrito};

% =================================================
% 6. BUZZER EN EL EXTREMO IZQUIERDO DEL CARRITO
% =================================================

% El buzzer conserva la altura de los micrófonos.
% Se ajusta la longitud del soporte para conectarlo
% con el cuerpo del carrito.

\draw[line width=0.85pt]
    (1.55,0.87) -- (1.55,1.47);

\node[
    circle,
    draw=orange!85!black,
    fill=orange!18,
    line width=1pt,
    minimum size=6mm,
    inner sep=0pt
] at (1.55,1.62) {};

\node[anchor=west] at (1.82,1.62)
    {Buzzer};

% =================================================
% 7. LIBERACION Y RECORRIDO
% =================================================

\node[anchor=south] at (1.55,2.32)
    {Liberación};

\draw[
    ->,
    line width=0.85pt
]
    (1.55,2.27) -- (1.55,1.96);

\draw[
    ->,
    line width=1pt
]
    (2.85,2.32) -- (8.15,2.32);

\node[
    fill=white,
    inner sep=3pt
]
    at (5.50,2.32)
    {Recorrido libre};

% =================================================
% 8. DISTANCIA ENTRE MICROFONOS: 200 cm
% =================================================

\draw[densely dashed,black!45]
    (0,-0.05) -- (0,-0.85);

\draw[densely dashed,black!45]
    (10,-0.05) -- (10,-0.85);

\node[
    fill=white,
    inner sep=2pt
]
    at (0,-0.58)
    {$0\,\mathrm{cm}$};

\node[
    fill=white,
    inner sep=2pt
]
    at (10,-0.58)
    {$200\,\mathrm{cm}$};

\draw[
    <->,
    line width=0.85pt
]
    (0,-0.85) -- (10,-0.85);

\node[
    fill=white,
    inner sep=3pt
]
    at (5,-1.10)
    {Separación entre micrófonos: $200\,\mathrm{cm}$};

% =================================================
% 9. LONGITUD UTIL: 180 cm
% =================================================

\draw[densely dashed,black!45]
    (0.70,-0.53) -- (0.70,-1.55);

\draw[densely dashed,black!45]
    (9.30,-0.53) -- (9.30,-1.55);

\draw[
    <->,
    line width=0.85pt
]
    (0.70,-1.43) -- (9.30,-1.43);

\node[
    fill=white,
    inner sep=3pt
]
    at (5,-1.70)
    {Longitud útil aproximada: $180\,\mathrm{cm}$};

% =================================================
% 10. PROTOCOLO EXPERIMENTAL
% =================================================

\node[
    draw=black!65,
    rounded corners=2pt,
    fill=black!2,
    line width=0.65pt,
    text width=3.65cm,
    align=left,
    inner sep=9pt,
    anchor=west
] at (11.05,1.45) {%
    \textbf{Protocolo previsto}\\[6pt]
    10 ensayos\\[3pt]
    Masa nominal: $200\,\mathrm{g}$\\[3pt]
    Adquisición: $10\,\mathrm{s}$
};

\end{tikzpicture}

    \caption{Esquema del montaje cinemático.
    El carro con la fuente acústica se desplaza sobre
    un riel entre dos micrófonos fijos. El movimiento
    se inicia mediante un empuje manual y el recorrido
    está limitado por topes mecánicos en sus extremos.}
    \label{fig:montaje_exp03}
\end{figure}

\subsubsection{Condiciones de adquisición}

Cada ensayo registra simultáneamente ambos canales
durante \SI{10}{\second}, a una tasa de
\SI{50}{kSPS} por canal después del diezmado
$M=10$. Esto equivale a 500\,000 muestras por
canal y ensayo. Se realizaron diez ensayos. El
buzzer permaneció activo durante toda la captura,
que comenzó antes del movimiento para conservar
un intervalo inicial de reposo. El carro se desplaza
mediante un empuje manual y puede recorrer el riel
en sentidos alternados; los topes rígidos delimitan
el movimiento.

La posición de la fuente acústica se obtuvo mediante
seguimiento óptico en Tracker. El seguimiento se
realiza sobre el buzzer. La secuencia original de
Tracker tiene una frecuencia de adquisición cercana
a \SI{30}{\hertz}; su interpolación a una base
temporal de \SI{100}{\hertz} permite compararla
con la salida cinemática de WALO, pero no incrementa
la resolución temporal ni la precisión espacial
intrínsecas del video. Los eventos de sincronización
identificados se utilizan como referencias candidatas
para alinear las series y no constituyen, por sí
solos, una medición validada de la latencia entre
sistemas.

\subsubsection{Estimación de velocidad y posición}

A partir de las señales registradas se estima
la frecuencia instantánea en cada canal mediante
el procesamiento temporal y espectral descrito
en el capítulo 1. Para expresar el desplazamiento
de frecuencia respecto al reposo, se utiliza una
referencia independiente por micrófono y por ensayo,
estimada en el intervalo inicial de reposo entre
0.50 y 0.80 s. Los valores medios de referencia
observados en la campaña fueron
\SI{2595.474}{\hertz} para MIC1
(desviación estándar \SI{1.988}{\hertz}) y
\SI{2596.610}{\hertz} para MIC2
(desviación estándar \SI{1.685}{\hertz}).
La variación de la frecuencia de reposo entre
ensayos y canales hace necesario conservar estas
referencias individuales y no imponer una única
frecuencia nominal a toda la campaña.

Con la convención de signo definida para cada
receptor, las velocidades radiales se calculan como

\begin{align}
    \hat{v}_1(t)&=
    -c(T)\frac{f_1(t)-f_{01}}{f_{01}},
    \label{eq:velocidad_mic1_exp03}\\
    \hat{v}_2(t)&=
    +c(T)\frac{f_2(t)-f_{02}}{f_{02}},
    \label{eq:velocidad_mic2_exp03}
\end{align}
donde $f_{01}$ y $f_{02}$ son las frecuencias de
referencia en reposo para MIC1 y MIC2,
respectivamente. Los signos opuestos corresponden
a la orientación enfrentada de los receptores y
a la convención de desplazamiento a lo largo del
riel. La combinación de ambos canales se obtiene
a partir de las estimaciones radiales individuales,
de acuerdo con la geometría del montaje.

La estimación depende de la estabilidad de las
referencias de reposo. Si la frecuencia del buzzer
deriva durante el ensayo, una variación propia de
la fuente puede confundirse con el desplazamiento
Doppler. Por ello, la frecuencia inicial de reposo
se estima por ensayo y canal, y la posible
no estacionariedad de esa referencia debe
considerarse al interpretar las velocidades.

Las estimaciones individuales representan la
componente radial de la velocidad bajo el modelo
Doppler de fuente móvil y receptor fijo. Su
interpretación requiere que la frecuencia de
referencia sea representativa del estado de reposo
y que las trayectorias acústicas dominantes sean
compatibles con la geometría idealizada. Además,
las señales recibidas corresponden a instantes de
emisión y propagación distintos; por ello, las
variaciones temporales de la frecuencia propia
del buzzer y las diferencias de propagación pueden
introducir discrepancias entre canales.

La posición puede estimarse mediante integración
temporal de la velocidad:

\begin{equation}
    \hat{x}(t)=x_0+
    \int_{t_0}^{t}\hat{v}_{\mathrm{diff}}(\tau)
    \,\mathrm{d}\tau,
    \label{eq:posicion_integrada_exp03}
\end{equation}

donde $x_0$ es la posición inicial de referencia.
La integración acumula los errores de velocidad,
por lo que cualquier sesgo sistemático puede
producir una desviación creciente en la posición
estimada.

La coordenada integrada se expresa respecto al
inicio de la captura y depende de la condición
inicial $x_0$ adoptada. En la comparación con
Tracker se debe mantener explícita la transformación
entre el origen de WALO, la geometría física del
riel y el origen de la coordenada óptica. La
corrección de deriva y el tratamiento de los
extremos del movimiento se evalúan a partir de
los registros; no se presupone que la integración
esté libre de deriva.

\subsubsection{Modelo dinámico de referencia}

El movimiento observado no corresponde a un
oscilador armónico con fuerza restauradora
elástica: el carro se impulsa manualmente y se
desplaza entre topes mecánicos. En consecuencia,
no se adopta un modelo masa--resorte ni se
interpretan los registros mediante una solución
exponencial sinusoidal.

Entre impactos, una descripción dinámica general
del movimiento puede escribirse como

\begin{equation}
    m\ddot{x}=F_f(v,\ldots),
    \label{eq:dinamica_libre_exp03}
\end{equation}

donde $F_f$ representa la fuerza resultante de
resistencia al movimiento, sin imponer aquí una
ley constitutiva específica. En los contactos
con los topes, una representación idealizada
del cambio de velocidad es

\begin{equation}
    v^{+}=-e\,v^{-},
    \label{eq:impacto_exp03}
\end{equation}

donde $v^{-}$ y $v^{+}$ son las velocidades
inmediatamente antes y después del impacto, y
$e$ es un coeficiente de restitución. Este
coeficiente no fue medido, por lo que la
expresión describe únicamente la estructura
idealizada de un impacto y no se utiliza para
ajustar parámetros.

La fuerza de resistencia puede incluir efectos
de fricción y pérdidas mecánicas, pero los datos
disponibles no justifican identificar por separado
un modelo de fricción seca o viscosa. Tampoco
se estiman la masa, un coeficiente de
amortiguamiento ni un coeficiente de restitución.
El análisis se centra en las magnitudes
cinemáticas medidas y en su comparación con
la referencia óptica.

\section{Tratamiento estadístico y métricas de error}
\label{sec:tratamiento_datos_error}

El análisis experimental requiere distinguir entre
la incertidumbre de las magnitudes de referencia,
la dispersión de las mediciones repetidas y el
error de los estimadores. Estas cantidades describen
aspectos diferentes del proceso de medición y no
deben utilizarse indistintamente.

\subsection{Estadística de las mediciones repetidas}

Para una magnitud $q$ medida $N$ veces, el valor
medio se calcula como

\begin{equation}
    \bar{q}=\frac{1}{N}\sum_{i=1}^{N}q_i.
    \label{eq:media_mediciones}
\end{equation}

La dispersión de las mediciones se caracteriza
mediante la desviación estándar muestral:

\begin{equation}
    s_q=
    \sqrt{
    \frac{1}{N-1}
    \sum_{i=1}^{N}(q_i-\bar{q})^2
    }.
    \label{eq:desviacion_estandar_mediciones}
\end{equation}

Cuando las mediciones pueden considerarse
independientes y representativas de una misma
condición, el error estándar de la media es

\begin{equation}
    \mathrm{SEM}=
    \frac{s_q}{\sqrt{N}}.
    \label{eq:sem_mediciones}
\end{equation}

La desviación estándar caracteriza la dispersión
de las observaciones individuales, mientras que
el SEM caracteriza la incertidumbre estadística
asociada a la estimación de la media bajo las
hipótesis indicadas. No son magnitudes
intercambiables.

\subsection{Error de estimación}

Para una magnitud lineal, el error individual
entre una estimación $q_i$ y su referencia
$q_{i,\mathrm{ref}}$ se define como

\begin{equation}
    e_i=q_i-q_{i,\mathrm{ref}}.
    \label{eq:error_individual}
\end{equation}

En el caso de magnitudes angulares, la diferencia
debe expresarse en el intervalo principal
$(-\pi,\pi]$ para evitar errores artificiales
al cruzar el límite de periodicidad. En la
expresión siguiente, los ángulos se expresan
en radianes. Se define

\begin{equation}
    \Delta\theta_i=
    \theta_i-\theta_{i,\mathrm{ref}},
\end{equation}

y el error angular reducido al intervalo principal como

\begin{equation}
    e_{\theta,i}=
    \operatorname{atan2}
    \left[
    \sin(\Delta\theta_i),
    \cos(\Delta\theta_i)
    \right].
    \label{eq:error_angular_principal}
\end{equation}

Las métricas angulares MAE, RMSE y sesgo se
calculan utilizando $e_{\theta,i}$ en lugar
de la diferencia angular sin envolver. Si se
reportan en grados, el error obtenido en radianes
se convierte antes de calcular o presentar
las métricas.

El error absoluto medio (MAE) y la raíz del
error cuadrático medio (RMSE) se calculan como

\begin{equation}
    \mathrm{MAE}=
    \frac{1}{N}\sum_{i=1}^{N}|e_i|,
    \label{eq:mae}
\end{equation}

\begin{equation}
    \mathrm{RMSE}=
    \sqrt{
    \frac{1}{N}\sum_{i=1}^{N}e_i^2
    },
    \label{eq:rmse}
\end{equation}

El sesgo medio se expresa mediante

\begin{equation}
    \mathrm{Bias}=
    \frac{1}{N}\sum_{i=1}^{N}e_i.
    \label{eq:bias}
\end{equation}

Estas métricas permiten evaluar la diferencia
entre las estimaciones y sus referencias.
No representan por sí mismas la incertidumbre
instrumental ni permiten identificar de forma
aislada el origen de los errores.

\subsection{Referencias geométricas e instrumentales}

En el experimento de distancia, la referencia
radial se determina mediante una cinta métrica
con resolución indicada de
\SI{1}{\milli\metre}. La incertidumbre efectiva
también depende de la alineación de los
transductores y de la determinación de sus
centros geométricos.

En el experimento angular, la referencia se
establece mediante las posiciones marcadas
respecto al centro del arreglo. La incertidumbre
angular de referencia se estima en aproximadamente
\SI{1}{\degree}. Este valor debe distinguirse
del error observado en los estimadores TDOA
y PDoA.

En el experimento cinemático, la geometría
mecánica del riel se describe mediante las
posiciones medidas de los límites efectivos
del recorrido. Esta geometría no equivale a
la coordenada óptica del buzzer ni al origen
temporal y espacial utilizado por WALO. La
posición obtenida mediante Tracker corresponde
al buzzer seguido en el video y debe tratarse
como una medición con resolución temporal y
espacial finita, no como un valor exacto. La
interpolación de los cuadros a una base de
\SI{100}{\hertz} facilita la comparación
numérica, pero no reduce la incertidumbre
asociada a la frecuencia original del video.
La masa del carro no es necesaria para la
estimación Doppler de velocidad y no se
dispone de una medición documentada de ella.

\subsection{Trazabilidad y conservación de los datos}

Los registros originales de cada experimento
se conservan como base del análisis. Las
transformaciones, filtros, estimadores y
criterios de selección deben poder relacionarse
con los datos de origen y con la configuración
de adquisición utilizada.

La evaluación de los resultados debe distinguir
entre datos de calibración, datos empleados
para ajustar los modelos y datos destinados
a validación independiente. Cuando no existe
una separación entre ajuste y evaluación,
los indicadores obtenidos deben identificarse
como resultados dentro de la muestra.

No se eliminan mediciones únicamente por
apartarse del comportamiento esperado. Cualquier
criterio de exclusión debe estar definido y
justificado por una condición instrumental
o experimental identificable. En ausencia
de dicho criterio, las anomalías se conservan
y se discuten como parte de la caracterización
del sistema.
